\subsection{Towards Optimal Transform}
\subsubsection{Karhunen-Lo\`eve Transform (KLT)}
The Karhunen--Lo\`eve Transform is the data-dependent orthogonal transform formed from covariance-matrix eigenvectors. It decorrelates transform coefficients and is optimal for energy compaction under its statistical model.
Assuming we have $X$: zero mean random vector of size $M$\\
Such that $R_X = E[XX^T]$, $R_X$ has M orthogonal eigenvectors $u_1 \cdots u_M$\\
Then KLT $\equiv$ an orthogonal Matrix $T_{KLT}$ where the rows are the eigenvectors:
$$T_{KLT} = \left[u_1 \ \vdots\ u_2\ \vdots\ \cdots\ \vdots\ u_M\right]^T$$
\textbf{Properties}:
\begin{itemize}
    \item $T_{KLT}^{-1} = T_{KLT}^T$
    \item Decorrelating transofrm, $E[y_iy_j] =\lambda_i \delta_{ij}$
    \item Best energy correlation (sparsity) $\forall N \in \{1 \cdots M\} : \sum E[y_i^2] \geq \sum E[z_i^2]$
\end{itemize}
3D intuition: take the spatial direction of point distribution and rotate it to the max,medium,min power axis. After transform we have sparsity and size reduction (Similar to ML PCA = Principal Component Analysis).\\
In the case of Gaussian data, this transform minimizes the geometric mean.\\
Example of usage: Multi/Hyper spectral imaging, where there is a large correlation across spectral bands.\\
KLT is not used typically in image codecs, due to:
\begin{itemize}
    \item Computational Complexity $\propto O(N^3)$ for eigenvectors, $O(N^2)$ for matrix multiplications
    \item Signaling overhead, $T_{KLT}$ sent apart
    \item Model Correctness: images are locally stationary
\end{itemize}
\subsection{Frequency Transforms}
\subsubsection{1D-DFT for Compression}
The Discrete Fourier Transform (DFT) represents a finite sequence as complex sinusoidal frequency components. Each coefficient measures amplitude and phase at one discrete frequency.
$$y[k] = \frac{1}{\sqrt{M}}\sr{M}{n=1}x[n]\exp\left(-j\frac{2\pi}{M}kn\right) \qquad \forall k = \{1 \cdots M\}$$
It is the same as a $T_{DFT}$ transform where:
$$T_{DFT} = \frac{1}{\sqrt{M}} \left[W_m^{i \cdot j}\right] \qquad \text{where} \begin{cases}
    i \in \{0 \cdots M-1\}\\
    j \in \{0 \cdots M-1\}\\
    W_m = \exp\left(-j\frac{2\pi}{M}\right)
\end{cases}$$
\subsubsection{2D-DFT for Compression}
$$y = Ax \qquad \text{where } A \text{ has size }N^2 \times N^2$$ 
\subsubsection{DFT Separable Implementation}
$$Y = \left(T X\right)T^T \qquad \qquad \text{with 2D basis func. }B_{kl}(n,m)=\frac{1}{N}\exp\left(j\frac{2\pi}{N}(kn+lm)\right)$$
First product $TX$ analyses the col. (V) frequency variation, while product with $T^T$ makes the row (H) analysis.\\
By sampling and periodizing the signal $\Rightarrow$ big jumps in the signal at period start/end $\Rightarrow$ \textbf{Frequency Leakage}.\\
In Frequency Leakage some energy leaks to the high frequencies (due to the big jumps = high variation).\\
Plotted DFT has a cross-like structure, usually it is plotted the $\log_{10}|DFT|$ version of the data\\
IDEA: why don't we put a mirror on the signal instead of periodizing, s.t. we reduce big jumps?
\subsubsection{DCT Transform Matrix}
The Discrete Cosine Transform (DCT) represents a finite real sequence using real cosine basis functions. Its symmetric boundary model reduces frequency leakage and usually compacts image energy better than a directly periodized DFT.
Steps to obtain it:
\begin{itemize}
    \item Create a mirrored-periodic version of $x \to x_{sym}$
    \item DFT of $x_{sym}$
    \item Apply frequency domain modulation (to obtain real coefficients)
    \item Keep only M coefficients
\end{itemize}
\subsubsection{1D-DCT Transform}
$$T_{DCT} = \forall n = \{0 \cdots M-1\} \to \begin{cases}
    \frac{1}{\sqrt{M}} \qquad\qquad\qquad\qquad k=0\\
    \sqrt{\frac{2}{M}}\cos\frac{\left(2n+1\right)k\pi}{2M} \qquad k>0
\end{cases}$$
\subsubsection{2D-DCT Transform: Block-Based DCT}
We just apply the DCT on both the horizontal and vertical axis.\\
Usually we apply an $n\times n$ block-based DCT transform, so with 64 vector of the basis.\\
The $n\times n$ matrix represents no variations with the UL corner, total horizontal variation on the UR corner, total vertical variation on the BL corner and a combination of total H/V variation in the BR corner.\\
Using the Block-Based DCT we get a very sparse signal, it is good because we apply the transform on blocks of the image: blocks are locally stationary, if there are more informations we use more bits, otherwise no.\\
How to decide the quantization step of each coefficient?
\begin{itemize}
    \item Adaptive Solutions: based on HS formula, or Greedy algorithm
    \item Fized Quantization Step: faster and used in JPEG
\end{itemize}
\subsection{JPEG}
JPEG is a lossy still-image coding standard based on block DCT, coefficient quantization, and lossless entropy coding. It defines a decodable bitstream syntax, guaranteeing interoperability while allowing different encoder implementations.\\
If we have 1 component $\to$ BW images, for color images $\to 3$ components: YCbCr.
\subsubsection{Encoding Strategy}
\begin{itemize}
    \item Take a block
    \item Subtrack the average value (128)
    \item Apply DCT on the blocks
    \item Quantization using Mid-Tread UQ $\to \tilde{C}_{ij} = \text{round}\left(\frac{C_{ij}}{q_{ij}}\right)$
    \item Zig-Zag Scan
    \item Entropy Coding (Huffman)
    \item Compressed Data in JPEG format
\end{itemize}
\begin{center}
\includegraphics[width=0.95\linewidth]{jpeg_pipeline.png}
\end{center}
\noindent Full JPEG encoding chain: RGB$\to$YCbCr conversion, chroma subsampling, $8\times8$ block partitioning, DCT, quantization (driven by the quantization table), zig-zag scan and entropy coding (Huffman/RLE) into the compressed JPEG stream.
\subsubsection{Quantization-step table}
The $q$ table is the quantization-step table: steps $\propto$ human sensitivity to the DCT components (64 byte table, $n \times n$ size).\\
Usually we use the $q^\star$ table, which has been made experimentally on scientific tests, bigger steps for lower sensitivity (BR corner), shorter steps for UL corner components.
\subsubsection{Scaling Factor and Quality}
JPEG quality $Q$ is an encoder control parameter, not an objective quality measurement. It scales quantization steps: larger $Q$ produces smaller steps, higher rate, and usually higher reconstructed quality; values are not directly comparable across different encoders.
The Quantization step table can be scaled by a factor $S_F$ depending on the quality we want:
$$S_F = \begin{cases}
    \frac{5000}{Q} \qquad\qquad 1 \leq Q \leq 50\\
    200 -2Q \qquad 50 < Q \leq 99\\
    1 \qquad \qquad\qquad Q = 100
\end{cases} \qquad \text{such that actual }q \leftarrow \frac{S_F}{100}q^\star$$
\subsubsection{Zig-Zag Coding}
Zig-zag scanning is an ordering of the $8\times8$ DCT coefficients from low toward high spatial frequencies. It tends to place long runs of quantized zeros at the end of the sequence, where run-length coding represents them cheaply.
We need a bitstream from the matrix. To encode the sparse matrix we use prediction + Huffman for the DC component, while the others are encoded using `run-length' coding.\\
Run-length coding: scanning sequentially the matrix and representing the value as $(i,j)$ where $i = $number of zeros before, and $j$ the actual value.\\
In this case the sequential scan is made with a zig-zag technique, final representation is:
$$\text{for the n-th block: }\qquad \left[ DC_{n} - DC_{n-1},\qquad (i,j) \to (\text{zig-zag}),\qquad \text{EOB: End of Block}\right]$$
